\subsection{Problem 5: Data Augmentation on ReducedMNIST}

This experiment evaluates how augmenting the reduced MNIST dataset affects classification performance for a fixed model architecture (LeNet-5). We measure three outcomes across a grid of base training-set sizes and additional synthetic images per class: (i) classification accuracy, (ii) F1 (macro), and (iii) wall-clock training time. The three base real dataset sizes considered are 350, 750, and 1000 examples (stratified across classes). Augmentation levels tested are 0 (no augmentation), 1000, 1500 and 2000 generated images per run. All experiments use the code in \texttt{Part2\_Problem5\_Augmentation/problem-5-final.ipynb} and a fixed random seed for reproducibility.

\subsubsection{Augmentation pipeline}
The augmentation pipeline applies one randomly selected transformation per image. The implemented transforms (and parameter ranges) are:
\begin{itemize}
    \item Rotate right: clockwise 5--30$^\circ$ (fill value 0)
    \item Rotate left: counter-clockwise 5--30$^\circ$ (fill value 0)
    \item Micro-rotate: small tilt $\pm3$--$5^\circ$
    \item Translate: random shift $\Delta x,\Delta y$ up to $\pm4$ pixels
    \item Rotate + Translate: 5--20$^\circ$ plus small shift up to $\pm3$ pixels
    \item Gaussian white noise: $\sigma$ in [0.02, 0.06] (image intensity scale)
    \item Salt-and-pepper: 2--5\% pixel flips to 0 or 1
    \item Rotate then noise: combination of rotation and Gaussian noise
\end{itemize}

These are implemented in the notebook's \texttt{AUGMENTATION\_FUNCS} set and are applied via a reproducible NumPy RNG per experiment.

\subsubsection{LeNet-5 model (as used)}
We use a compact LeNet-5 variant defined in the notebook. Layer summary (shapes assume input $1\times28\times28$):
\begin{itemize}
        \item Conv1: Conv2d(1, 6, kernel=5, padding=2) -- activation: Tanh; AvgPool2d(2)
        \item Conv2: Conv2d(6, 16, kernel=5) -- activation: Tanh; AvgPool2d(2)
        \item Classifier: Flatten $\rightarrow$ Linear(16$\times$5$\times$5, 120) -- Tanh $\rightarrow$ Linear(120,84) -- Tanh $\rightarrow$ Linear(84,10)
        \item Loss: CrossEntropy; Optimizer: Adam; LR schedule: CosineAnnealingLR
\end{itemize}

\subsubsection{Evaluation metrics}
We report two standard metrics for multiclass classification plus wall-clock training time.
\begin{itemize}
    \item \textbf{Accuracy}: fraction of correctly predicted examples,
    $$\mathrm{Accuracy} = \frac{1}{N}\sum_{i=1}^N \mathbf{1}\!\left\{\hat{y}_i = y_i\right\},$$
    where $\hat{y}_i$ is the predicted class, $y_i$ the ground-truth, and $N$ the number of examples.
    \item \textbf{F1 (macro)}: per-class F1 averaged across classes. Let for class $c$ precision $P_c$ and recall $R_c$ be defined as usual; the per-class F1 is
    $$F_{1,c} = \frac{2P_c R_c}{P_c + R_c},$$
    and the macro score is
    $$F1_{\text{macro}} = \frac{1}{C} \sum_{c=1}^C F_{1,c},$$
    where $C$ is the number of classes (10 here). Macro F1 weights all classes equally regardless of support.
\end{itemize}

\subsubsection{Results}
Below are the three result tables produced from the notebook outputs: Accuracy (\%), F1 (macro) and Time (s). Replace the placeholder cells with the numeric outputs from the notebook \texttt{pivot\_acc}, \texttt{pivot\_f1}, and the timing values saved for each experiment.

\begin{table}[h]
        \centering
        \caption{Accuracy (\%) across augmentation levels and base dataset sizes}
        \label{tab:aug_accuracy}
        \begin{tabular}{|l|c|c|c|}
            \hline
            \textbf{n\_aug / n\_real} & \textbf{350 Real} & \textbf{750 Real} & \textbf{1000 Real} \\
            \hline
            0    & 82.5 & 92.0 & 91.5 \\
            1000 & 87.5 & 92.0 & 92.5 \\
            1500 & 88.0 & 93.0& 93.0 \\
            2000 & 88.5 & 94.0 & 94.5 \\
            \hline
        \end{tabular}
\end{table}

\begin{table}[h]
        \centering
        \caption{F1 (macro) across augmentation levels and base dataset sizes}
        \label{tab:aug_f1}
        \begin{tabular}{|l|c|c|c|}
            \hline
            \textbf{n\_aug / n\_real} & \textbf{350 Real} & \textbf{750 Real} & \textbf{1000 Real} \\
            \hline
            0    & 0.8209 & 0.9202 & 0.9142 \\
            1000 & 0.8740 & 0.9199 & 0.9242 \\
            1500 & 0.8790 & 0.9310 & 0.9296 \\
            2000 & 0.8846 & 0.9395 & 0.9446 \\
            \hline
        \end{tabular}
\end{table}

\begin{table}[h]
        \centering
        \caption{Training time (s) across augmentation levels and base dataset sizes}
        \label{tab:aug_time}
        \begin{tabular}{|l|c|c|c|}
            \hline
            \textbf{n\_aug / n\_real} & \textbf{350 Real} & \textbf{750 Real} & \textbf{1000 Real} \\
            \hline
            0    & 1.5 & 0.5 & 0.6 \\
            1000 & 1.3 & 1.4 & 1.6 \\
            1500 & 1.7 & 1.9 & 2.1 \\
            2000 & 2.2 & 2.4 & 2.6 \\
            \hline
        \end{tabular}
\end{table}

\subsubsection{Discussion}
The empirical results in Tables~\ref{tab:aug_accuracy} and~\ref{tab:aug_f1} show consistent improvements from augmentation, with the largest relative gains at the smallest base size (350 examples). Quantitative observations:
\begin{itemize}
    \item For 350 real examples accuracy increases from 82.5\% (no augmentation) to 88.5\% at 2000 augmented images (absolute gain 6.0\%). F1 (macro) increases from 0.8209 to 0.8846 (gain 0.0637).
    \item For 750 real examples accuracy increases from 92.0\% to 94.0\% (gain 2.0\%); F1 rises from 0.9202 to 0.9395 (gain 0.0193).
    \item For 1000 real examples accuracy increases from 91.5\% to 94.5\% (gain 3.0\%); F1 rises from 0.9142 to 0.9446 (gain 0.0304).
\end{itemize}

Training time (Table~\ref{tab:aug_time}) grows with the amount of augmentation. Example increases from no augmentation to 2000 augmented images:
\begin{itemize}
    \item 350 real: 1.5 s -> 2.2 s (+0.7 s)
    \item 750 real: 0.5 s -> 2.4 s (+1.9 s)
    \item 1000 real: 0.6 s -> 2.6 s (+2.0 s)
\end{itemize}

Taken together, augmentation provides clear accuracy/F1 benefits, especially when the base dataset is small. However, marginal gains shrink as the real data pool grows. The 1500-2000 augmented range yields only small improvements over 1000--1500 in several cases, suggesting diminishing returns beyond ~1500 synthetic images for this model/dataset combination.

\subsubsection{Conclusions}
The results indicate that controlled augmentation reliably improves classification performance for the LeNet-5 model on ReducedMNIST. Key takeaways:
\begin{itemize}
    \item Augmentation yields the largest absolute benefit at the smallest base size (350): +6.0\% accuracy and +0.0637 F1 when moving to 2000 augmented images.
    \item For medium/large base sizes (750, 1000) the gains are smaller (2--3\% accuracy), implying that augmentation is most valuable when real data is scarce.
    \item Training time increases with augmentation; the 1500-2000 range yields limited extra performance versus the 1000--1500 range, so 1500 is a reasonable compromise for this setup.
\end{itemize}

Overall recommendation: apply augmentation (1000--1500 synthetic images) when limited to a few hundred real samples, and re-evaluate marginal gains before scaling beyond 1500 synthetic samples per experiment.

% End of section
